\section*{第 \ref{ch:b2:poynting-vector} 章——电磁能量与坡印廷矢量}

\begin{solution}{exo:b2:poynting-vector:1}
$E_0 = \sqrt{2I/\varepsilon_0c} = \qty{870}{V/m}$，$B_0 = E_0/c = \qty{2.9}{\micro T}$；$u = I/c
= \qty{3.3}{\micro J/m^3}$；$P = I/c = \qty{3.3}{\micro Pa}$（全吸收），\qty{6.7}{\micro Pa}
（镜面）；屋顶上为 \qty{0.33}{mN}；$I/hf = 1000/3.2 \times 10^{-19} = \qty{3e21}{m^{-2}.s^{-1}}$。
\end{solution}

\begin{solution}{exo:b2:poynting-vector:2}
面积 \qty{1.8}{mm^2}：$I = \qty{2.8}{kW/m^2}$，$E_0 = \qty{1.5}{kV/m}$，$B_0 = \qty{4.9}{\micro T}$。
聚焦至 \qty{5}{\micro m}：$I = \qty{2.5e8}{W/m^2}$，$E_0 = \qty{4.4e5}{V/m}$。拍瓦级激光：
$I = \qty{1e26}{W/m^2}$，$E_0 = \qty{2.7e14}{V/m}$，为原子场的五百倍——
原子在一个周期内即被撕裂。
\end{solution}

\begin{solution}{exo:b2:poynting-vector:3}
$\tfrac12\varepsilon_0E^2$：\qty{40}{J/m^3}，\qty{4.4e4}{J/m^3}。$B^2/2\mu_0$：\qty{4}{kJ/m^3}，
\qty{0.9}{MJ/m^3}，\qty{8e8}{J/m^3}。磁共振成像仪的孔径内储存有 \qty{0.9}{MJ} 的能量——相当于 \qty{26}{mL}
汽油的能量。
\end{solution}

\begin{solution}{exo:b2:poynting-vector:4}
$F = 2IA/c = \qty{0.09}{N}$；$a = \qty{4.5e-4}{m/s^2}$；一个月内速度增加 \qty{1.2}{km/s}；
太阳的拉力为 \qty{1.2}{N}：帆无法与引力抗衡，
而是在轨道上逆着引力斜向航行。
\end{solution}

\begin{solution}{exo:b2:poynting-vector:5}
(a) $j = \qty{3.2e6}{A/m^2}$，$E = j/\gamma = \qty{0.053}{V/m}$；$B = \mu_0I/2\pi a = \qty{2}{mT}$；
$\Pi = EB/\mu_0 = \qty{84}{W/m^2}$，沿径向向内。(b) $84 \times 2\pi a = \qty{0.53}{W}$
每米；$R = 1/\gamma\pi a^2 = \qty{5.3}{m\ohm/m}$，$RI^2 = \qty{0.53}{W}$。(c) $\Pi(r)
= EIr/2\pi a^2$，故 $2\pi r\Pi = EIr^2/a^2$，且壳层接收的功率为 $\dd(2\pi r\Pi)
= 2EIr\,\dd r/a^2 = (j^2/\gamma)2\pi r\,\dd r$。(d) 电池表面的电荷
在整个电路中、导线外部建立电场；
该处 $\vect E\wedge\vect B$ 将能量从电池输送至负载。
\end{solution}

\begin{solution}{exo:b2:poynting-vector:6}
(a) $E = Q/\varepsilon_0\pi R^2$ 沿轴向，$B = \mu_0Ir/2\pi R^2$ 沿方位角方向，$\vect\Pi = -\varepsilon_0
(r/2)E\dot E\,\vect e_r$。(b) $\varepsilon_0(R/2)E\dot E\cdot2\pi Rd = \dd(\tfrac12\varepsilon_0E^2\pi R^2d)/
\dd t$。(c) $\dot E = I/\varepsilon_0\pi R^2 = \qty{1.4e13}{V.m^{-1}.s^{-1}}$；$\Pi(R) = \qty{3.2e6}{W/m^2}$；
乘以 $2\pi Rd = \qty{3.1e-4}{m^2}$：\qty{1.0}{kW} $= UI$，其中 $U = Ed = \qty{1}{kV}$。(d)
是的：小了一个 $(\omega R/c)^2$ 类的因子（见\cref{ch:b2:maxwell-equations}）。
\end{solution}

\begin{solution}{exo:b2:poynting-vector:7}
(a) $\vect E = -(r/2)\dot B\,\vect e_\theta$，$\vect\Pi = -(rB\dot B/2\mu_0)\vect e_r$：能流向内，
当 $B$ 增大时。(b) $(RB\dot B/2\mu_0)2\pi R = \dd(\pi R^2B^2/2\mu_0)/\dd t$。(c)
在外部 $\vect B = \vect 0$，故 $\vect\Pi = \vect 0$：能量恰好在
绕组处注入，发电机驱动电流逆着
感应场（该处 $\vect j\cdot\vect E < 0$）；紧贴薄片内侧，
$\vect\Pi$ 已是 (a) 中的向内能流。
\end{solution}

\begin{solution}{exo:b2:poynting-vector:8}
(a) $E = U/r\ln(b/a)$ 沿径向，$B = \mu_0I/2\pi r$ 沿方位角方向。(b) $\vect\Pi = (UI/
2\pi r^2\ln(b/a))\vect e_z$，沿纤芯电流方向。(c) $\int_a^b\Pi\,2\pi r\,\dd r =
UI$。(d) $E$ 获得一个轴向分量 $j/\gamma$；$\vect\Pi$ 向纤芯倾斜，
其流入每米纤芯的径向通量为 $RI^2$（每米）。
\end{solution}

\begin{solution}{exo:b2:poynting-vector:9}
(a) $F_{\text{rad}} = I_0(D_0/D)^2\pi a^2/c$，$F_g = GM\cdot\tfrac43\pi a^3\rho/D^2$：比值为
$3I_0D_0^2/4cGM\rho a$，与 $D$ 无关。(b) 比值 $> 1$ 的条件是 $a < 3I_0D_0^2/
4cGM\rho = \qty{0.3}{\micro m}$。(c) 细尘埃被推离太阳，
同时保持彗星的轨道速率，因此尘埃尾背离太阳，
并弯曲；受太阳风驱动的离子尾则是直的。
\end{solution}

\begin{solution}{exo:b2:poynting-vector:10}
(a) $\dot E < 0$：$\vect\Pi$ 在边缘处指向外。(b) $-\dd W_e/\dd t$
离开电容器，$RI^2$ 进入导线；在准静态
条件下两者相等——二者之间不积累能量。(c) $W_e =
Q^2/2C$，$\dd W_e/\dd t = (Q/C)\dd Q/\dd t = -UI = -RI^2$。(d) $\tfrac12LI^2$
在两端（$I = 0$）处为零，
对小回路而言其中间值可忽略。
\end{solution}

\begin{solution}{exo:b2:poynting-vector:11}
(a) $\vect\Pi = (EB/\mu_0)\vect e_y$，均匀：$\operatorname{div}\vect\Pi = 0$。(b) 零。
(c) 沿 $y$ 方向穿过间隙，并经
由极板与磁铁边缘处的边缘场闭合。(d) 静止于静态场中的物体
不携带电流：$\vect j\cdot\vect E = 0$，无发热。
环流的 $\vect\Pi$ 只是记账量，并非沉积能量的流动。
\end{solution}

\begin{solution}{exo:b2:poynting-vector:12}
法拉第：$\partial_xE_y = -2E_0k\sin kx\cos\omega t$ 等于 $-\partial_tB_z$，因为
$\omega = kc$。(a)
$\vect\Pi = (E_yB_z/\mu_0)\vect e_x = -(E_0^2/\mu_0c)\sin2kx\sin2\omega t\,\vect e_x$：平均值为
零。(b) $u = 2\varepsilon_0E_0^2(\cos^2kx\cos^2\omega t + \sin^2kx\sin^2\omega t)$：电场能量
位于 $x = 0, \lambda/2, \dots$ 处，磁场能量位于 $\lambda/4, \dots$ 处，
两者之间每周期往复转移两次。(c) 动量：每列波的 $I = \varepsilon_0cE_0^2/2$，
反射波使其动量反向：$2I/c = \varepsilon_0E_0^2$。拉普拉斯力：在 $x =
0$ 处，紧贴外侧 $B = (2E_0/c)\sin\omega t$，内侧为 $0$，故 $j_s = B/\mu_0$；
薄片单位面积所受的力为 $\tfrac12j_sB = B^2/2\mu_0 = 2\varepsilon_0E_0^2\sin^2
\omega t$。(d) 其平均值为 $\varepsilon_0E_0^2 = 2I/c$。
\end{solution}

\begin{solution}{pb:b2:poynting-vector:1}
\textbf{1.} $E = U/r\ln(b/a)$：在 $r = a$ 处为 \qty{1.8}{MV/m}；$B = \mu_0I/2\pi r$：
在 $a$ 处为 \qty{10}{mT}。

\textbf{2.} $\vect\Pi = (UI/2\pi r^2\ln(b/a))\vect e_z$；$\int_a^b\Pi\,2\pi r\,\dd r = UI =
\qty{10}{MW}$。

\textbf{3.} $\tfrac12\varepsilon_0\varepsilon_rE^2 = \qty{34}{J/m^3}$，$B^2/2\mu_0 = \qty{40}{J/m^3}$：
量级相当。

\textbf{4.} $\Pi(a) = EB/\mu_0 = \qty{1.45e10}{W/m^2}$，$u = \qty{74}{J/m^3}$：$\Pi/u =
\qty{2.0e8}{m/s} = c/\sqrt{\varepsilon_r}$。

\textbf{5.} $E_z = I/\gamma\pi a^2 = \qty{0.027}{V/m}$; $\Pi_r = E_zB/\mu_0 = \qty{210}{W/m^2}$
指向内侧；$\times2\pi a = \qty{13}{W/m}$；$R = 1/\gamma\pi a^2 = \qty{53}{\micro\ohm/m}$，$RI^2 =
\qty{13}{W/m}$.

\textbf{6.} $\Pi_r/\Pi_z = E_z/E_r = 1.5 \times 10^{-8}$：几乎全部能量
沿导线传输，仅极少量渗入铜中。

\textbf{7.} 损失 \qty{133}{kW}，占 \qty{10}{MW} 的 $1.3\%$；在 \qty{200}{kV} 下为 \qty{1.3}{kW}，
$0.013\%$.

\textbf{8.} $\int_a^b(B^2/2\mu_0)2\pi r\,\dd r = (\mu_0I^2/4\pi)\ln(b/a) = \qty{27}{mJ/m} =
\tfrac12\Lambda I^2$，其中 $\Lambda = \mu_0\ln(b/a)/2\pi = \qty{0.22}{\micro H/m}$；电场情形：
$\tfrac12\Gamma U^2$，其中 $\Gamma = 2\pi\varepsilon_0\varepsilon_r/\ln(b/a) = \qty{117}{pF/m}$：\qty{23}{mJ/m}。

\textbf{9.} $E_0 = \qty{870}{V/m}$, $B_0 = \qty{2.9}{\micro T}$, $u = \qty{3.3}{\micro J/m^3}$,
\qty{3e21}{m^{-2}.s^{-1}}.

\textbf{10.} 电功率 \qty{320}{W}；热功率 \qty{1280}{W}，经 \qty{3.2}{m^2}
表面散出：$\Delta T = 1280/32 = \qty{40}{K}$。

\textbf{11.} $P = I/c = \qty{3.3}{\micro Pa}$；\qty{5.3}{\micro N}，与
\qty{180}{N}.

\textbf{12.} $320 \times 6 = \qty{1.9}{kWh/day}$，\qty{700}{kWh/year}；五块太阳能板，面积
\qty{8}{m^2}.

\textbf{13.} $30 \times 10.7 = \qty{320}{W}$；功率在两导线之间的场中流动，
即 $\vect E$ 跨接 \qty{30}{V} 电压，$\vect B$ 绕着电流——
而不在铜中。

\textbf{14.} $4\pi D^2I = \qty{3.9e26}{W}$; $P/c^2 = \qty{4.3e9}{kg/s}$.

\textbf{15.} $5 \times 10^{21} \times 1.6 = \qty{5e21}{}$ 每秒光子数。

\textbf{16.} $I\pi R_E^2/c = \qty{5.8e8}{N}$：相当于重力的 $10^{-14}$。

\textbf{17.} $\lambda = \qty{12.2}{cm}$；\qty{30}{cm} 的谐振腔约为几个波长
——是具有若干模式的谐振腔，而非自由光束；金属壁反射
电磁波（趋肤深度为微米量级），将其约束在腔内。

\textbf{18.} $0.25 \times 4180 \times 60 = \qty{63}{kJ}$，功率为 \qty{640}{W} 时需 \qty{98}{s}；
$640/2.5 \times 10^{-4} = \qty{2.6}{MW/m^3}$.

\textbf{19.} $I = 640/5 \times 10^{-3} = \qty{1.3e5}{W/m^2}$，$E_0 = \qty{10}{kV/m}$；热点
相距 $\lambda/2 = \qty{6}{cm}$——因此需要转盘。

\textbf{20.} $W = QP/\omega = 10^3 \times 800/1.54 \times 10^{10} = \qty{52}{\micro J}$; $u =
\qty{1.7}{mJ/m^3}$；$E_0 = \sqrt{2u/\varepsilon_0} = \qty{20}{kV/m}$——比击穿场强
低一百倍。

\textbf{21.} $\tfrac12\omega\varepsilon_0\varepsilon''E_0^2 = \qty{2.6e6}{W/m^3}$: $E_0 = \qty{2}{kV/m}$
水中——与腔场同量级，因水的介电常数而减弱，
结果相符。

\textbf{22.} $hf = \qty{1.6e-24}{J} = \qty{1e-5}{eV}$；每秒 $4 \times 10^{26}$ 个光子；
每个光子的能量比电离所需的少约一百万倍。

\textbf{23.} $\sim u = \qty{2}{mPa}$：否。

\textbf{24.} 若无物质吸收它，$\vect\Pi$ 的通量
别无去处，只能流回磁控管，致使其过热。

\textbf{25.} 传输线：电介质中，$\vect\Pi$ 沿轴向，\qty{10}{MW}，末端接
负载。太阳能板：光束中，$\vect\Pi$ 沿太阳光线，\qty{1.6}{kW}
输入，\qty{320}{W} 输出，半导体。烤箱：腔场中，$\vect\Pi$
环绕流动，\qty{800}{W}，水的偶极子。
\end{solution}
